\LevelZero{Permission}\label{ch:permission}
\LevelOne{Definition}
A \ttt{permission} is used to decide whether \ttt{a player} can do something or not.


\LevelOne{Types}
To make the discussion clearer, we define the types of \ttt{permission} as follows:
\begin{description}
    \item [level permission] A permission level is a non-negative number used in vanilla Minecraft.
    The higher number means the higher authority.
    \item[string permission] Usually, a \ttt{string permission} is introduced by a \ttt{permission plugin}, such as \ttt{luckperms} mod.
\end{description}
\clearpage

\LevelOne{What is the permission system used by fuji?}
\LevelTwo{Explanation}
Fuji use the mojang's \ttt{vanilla Minecraft permission system}, which is based on \ttt{level permission}.\\
\\
As a convention, most of the commands registered by fuji, requires \ttt{level permission} to be \ttt{0} to use.
Only a few of the commands require the \ttt{level permission} to be \ttt{4} to use.

\begin{tips}{Modify the default required level permission to be 4 for all fuji commands}
   You may want to prohibit the player from using any commands registered by fuji.
   And then assign the proper permission for each command one by one.
   In this case, you can set the \str{all\_commands\_require\_level\_4\_permission\_to\_use\_by\_default} in \mainconfig file to be \tbf{true}.
\end{tips}

\LevelTwo{Set a string permission for a command}
By default, fuji only use the \ttt{level permission} as \ttt{the requirement of} a command.
However, if you want to use \ttt{string permission} for a \ttt{command}, you can use \tbf{command\_permission} module, to provide luckperms permissions for each command.

\begin{tips}{Query all registered string permissions.}
    Issue \cmd{/fuji inspect permissions-and-metas} command, to see all the registered permissions.
\end{tips}



\LevelOne{Reference}
\begin{enumerate}
    \item \url{https://minecraft.fandom.com/wiki/Permission\_level}
    \item \url{https://luckperms.net/wiki/Home}
\end{enumerate}

